%% kapa: exam=1 round=28 year=2017 period=1 form=A subject=real_estate subject_ko=부동산학원론 seq=27 number=107
%% source: Q-Net 감정평가사 기출문제 (cst00309) artlSeq=5199603 file=extracted/1008/5199603/2017년 제28회 감정평가사 1차시험 문제지/1교시 A형.pdf
%% conversion: kordoc 4.24.51 -> tools/gen1.py (auto); review=proofread-sonnet; flags=none
\begin{question}{107}
부동산개발과 시장 분석에 관한 설명으로 옳지 않은 것은? (단, 주어진 조건에 한함)
\begin{choices}
  \item 부동산 개발과정에서 시장분석의 목적은 개발과 관련된 의사결정을 하기 위하여 부동산의 특성상 용도별, 지역별로 각각의 수요와 공급에 미치는 요인들과 수요와 공급의 상호 관계가 개발사업에 어떠한 영향을 미치는가를 조사ㆍ분석하는 것이다.
  \item 시장성분석은 현재와 미래의 대상 부동산에 대한 수요ㆍ공급 분석을 통해 흡수율 분석과 시장에서 분양될 수 있는 가격, 적정개발 규모 등의 예측을 한다.
  \item 지역경제분석은 지역의 경제활동, 지역인구와 소득 등 대상지역시장 전체에 대한 총량적 지표를 분석한다.
  \item 부동산 개발과정의 시장분석은 속성상 지리적ㆍ공간적 범위에 국한되지 않으며, 대상 개발사업의 경쟁력 분석에 한한다.
  \item 경제성분석은 구체적으로 개발사업의 수익성 여부 등을 평가한다.
\end{choices}
\end{question}
\answer{4}
